% from ega4-8-fr.tex
\oldpage[IV-3]{22}
et $u_\nu^*(f_\nu\circ g_\nu)=1_{\mathcal G}$~; compte tenu de ce que
$\mathcal F_\nu$ et $\mathcal G_\nu$ sont de type fini, on en conclut par
(8.5.2, (i)) qu'il existe $\mu\geq\nu$ tel que l'on ait
$g_\mu\circ f_\mu=1_{\mathcal F_\mu}$ et
$f_\mu\circ g_\mu=1_{\mathcal G_\mu}$, d'où le corollaire.
\end{proof}

\begin{corollary}[8.5.2.5]
\label{IV.8.5.2.5-fr}
Supposons $S_0$ quasi-compact et quasi-séparé. Supposons que
$\mathcal F_\lambda$, $\mathcal G_\lambda$ soient des
$\mathcal O_{S_\lambda}$-Modules quasi-cohérents de présentation finie. Pour
que $\mathcal F$ et $\mathcal G$ soient isomorphes, il faut et il suffit qu'il
existe $\mu\geq\lambda$ tel que $\mathcal F_\mu$ et $\mathcal G_\mu$ soient
isomorphes. De plus, pour tout isomorphisme
$f:\mathcal F\xrightarrow{\sim}\mathcal G$, il existe un $\nu\geq\mu$ et un
isomorphisme $f_\nu:\mathcal F_\nu\xrightarrow{\sim}\mathcal G_\nu$ tels que
$f=u_\nu^*(f_\nu)$.
\end{corollary}

\begin{proof}
Cela résulte de (8.5.2.4) et de (8.5.2, (i)) puisque tout homomorphisme
$f:\mathcal F\to\mathcal G$ est de la forme $u_\mu^*(f_\mu)$ pour un
$\mu\geq\lambda$ et un homomorphisme
$f_\mu:\mathcal F_\mu\to\mathcal G_\mu$.
\end{proof}

\par\medskip\noindent\textup{(ii)} Considérons encore en premier lieu le cas où
$S_0=\operatorname{Spec}(A_0)$ est
affine. Alors l'assertion équivaut au lemme (5.13.7.1).

Dans le cas général, en partant d'un recouvrement ouvert affine fini $(U_i)$
de $S_0$, on déduit de (5.13.7.1) que pour tout $i$, il existe un indice
$\lambda(i)$ et un $\mathcal O_{U_{i,\lambda(i)}}$-Module quasi-cohérent de
présentation finie $\mathcal F_i$ tels que
$u_{\lambda(i)}^*(\mathcal F_i)=\mathcal F|V_i$ (avec les notations de (i)). En
outre, comme $L$ est filtrant, on peut supposer que tous les $\lambda_i$ sont
égaux à un même $\lambda$. Comme
$U_{ij\lambda}=U_{i\lambda}\cap U_{j\lambda}$ est quasi-compact et
quasi-séparé (1.2.7), il résulte de (8.5.2.5) que pour tout couple $(i,j)$, il
existe un indice $\lambda_{ij}\geq\lambda$ et un isomorphisme
\[
  \theta_{ij}:u_{\lambda,\lambda_{ij}}^*(\mathcal F_i|U_{ij\lambda})
  \xrightarrow{\sim}
  u_{\lambda,\lambda_{ij}}^*(\mathcal F_j|U_{ij\lambda})
\]
tel que $u_{\lambda_{ij}}^*(\theta_{ij})$ soit l'automorphisme identique de
$\mathcal F|(V_i\cap V_j)$~; on peut encore supposer tous les
$\lambda_{ij}$ égaux à $\lambda$. Changeant les notations, on peut donc
supposer qu'il existe pour tout couple $(i,j)$ un isomorphisme
\[
  \theta_{ij}:\mathcal F_i|U_{ij\lambda}
  \xrightarrow{\sim}\mathcal F_j|U_{ij\lambda},
\]
tels que $u_\lambda^*(\theta_{ij})$ soit l'automorphisme identique de
$\mathcal F|(V_i\cap V_j)$. Enfin, pour trois indices quelconques $i,j,k$, si
l'on pose
$U_{ijk,\lambda}=U_{i\lambda}\cap U_{j\lambda}\cap U_{k\lambda}$,
$U_{ijk,\lambda}$ est quasi-compact, et si $\theta'_{ij}$, $\theta'_{jk}$ et
$\theta'_{ik}$ désignent les restrictions de $\theta_{ij}$, $\theta_{jk}$ et
$\theta_{ik}$ à $U_{ijk,\lambda}$, on a
\[
  u_\lambda^*(\theta'_{ij}\circ\theta'_{jk})=u_\lambda^*(\theta'_{ik}).
\]
Il y a donc, en vertu de (i), un indice $\mu\geq\lambda$ tel que l'on ait
\[
  u_{\lambda\mu}^*(\theta'_{ik})
  =u_{\lambda\mu}^*(\theta'_{ij}\circ\theta'_{jk}),
\]
donc les isomorphismes
\[
  u_{\lambda\mu}^*(\theta'_{ij}):
  u_{\lambda\mu}^*(\mathcal F_i)|U_{ij\mu}
  \xrightarrow{\sim}
  u_{\lambda\mu}^*(\mathcal F_j)|U_{ij\mu}
\]
vérifient la condition de recollement, et définissent par suite sur $S_\mu$ un
$\mathcal O_{S_\mu}$-Module quasi-cohérent $\mathcal F_\mu$ de présentation
finie $(0_{\mathrm I},3.3.1)$ tel que $\mathcal F$ soit isomorphe à
$u_\mu^*(\mathcal F_\mu)$.

\begin{env}[8.5.3]
\label{IV.8.5.3-fr}
Le résultat de (8.5.2) s'exprime encore en disant que si $S_0$ est
quasi-compact et quasi-séparé, la catégorie des $\mathcal O_S$-Modules
quasi-cohérents de présentation finie est déterminée à équivalence près par la
donnée des catégories des $\mathcal O_{S_\lambda}$-Modules quasi-cohérents de
présentation finie, des foncteurs $u_{\lambda\mu}^*$ entre ces catégories, et
des isomorphismes de transition
$u_{\mu\nu}^*\circ u_{\lambda\mu}^*\xrightarrow{\sim}u_{\lambda\nu}^*$. De façon
imagée, on peut dire que la donnée d'un $\mathcal O_S$-Module quasi-cohérent de
présentation finie $\mathcal F$ équivaut «~fonctoriellement~» à la donnée d'un
$\mathcal O_{S_\lambda}$-Module de présentation finie $\mathcal F'_\lambda$
pour $\lambda$ grand et que, si un $\mathcal O_{S_\mu}$-Module quasi-cohérent de
présentation finie $\mathcal F''_\mu$ a aussi $\mathcal F$ comme image
réciproque, alors $\mathcal F'_\lambda$ et $\mathcal F''_\mu$ ont même image
réciproque dans un $S_\nu$ convenable ($\nu\geq\lambda$, $\nu\geq\mu$).

Nous allons interpréter de ce point de vue diverses notions liées aux
$\mathcal O_S$-Modules quasi-cohérents.
\end{env}

